本文针对红外反射光谱测量半导体外延层厚度的问题，建立由双光束正演、参数反演到多光束影响分析逐步递进的光学测厚方法。针对SiC与Si在不同吸收条件下的光学响应，综合采用Fresnel干涉、有界非线性最小二乘和Airy多光束模型求取外延层厚度，并分析模型近似对反演结果的影响。

\textbf{针对问题一，}建立\textbf{Sellmeier--Drude--Fresnel双光束干涉模型}。由空气--外延层--衬底三层光路推导光程差和传播相位，并以Sellmeier色散与Drude自由载流子响应描述SiC折射率。将表面反射光和第一束内部反射光按Fresnel振幅系数相干叠加，得到非偏振理论反射率。弱色散条件下局部条纹间距与厚度近似成反比。该关系为后续利用连续干涉光谱反演外延层厚度提供了理论依据。

\textbf{针对问题二，}建立\textbf{SiC双角度有界非线性最小二乘厚度反演模型}。根据弱吸收条件选取$2500$--$3300\,\mathrm{cm}^{-1}$有效波段，并分别处理$10^\circ$和$15^\circ$两组实测光谱。以理论反射率与实测反射率的残差平方和为目标函数，在厚度及相关光学参数的物理范围内进行约束反演。采用“低维盆地搜索--完整参数精修”的两阶段多初值算法识别不同干涉级次对应的局部极小解。两个角度分别得到$7.8566\,\mu\mathrm m$和$7.6230\,\mu\mathrm m$，综合得到SiC外延层厚度\textbf{$7.7398\,\mu\mathrm m$}，两组拟合的$R^2$均高于$0.95$、MAPE均低于$0.30\%$。两角度厚度相对差异为$3.02\%$，且最优厚度解在低维搜索与完整参数精修中均稳定出现，表明反演结果具有较好的跨角度一致性。

\textbf{针对问题三，}建立\textbf{双振子--Drude复折射率与Airy多光束干涉模型}。由高阶反射场递推关系分析界面反射、传播衰减和相干叠加条件，并以复折射率同时描述硅的色散与吸收。对无限多束相干反射求和得到Airy反射率，据此建立硅光谱的四参数有界反演模型。采用多初值非线性最小二乘求解，并通过双光束截断误差传播和Jacobian--SVD分析多光束影响及参数可辨识性。两个角度分别得到$3.2474\,\mu\mathrm m$和$3.1875\,\mu\mathrm m$，综合得到硅外延层厚度\textbf{$3.2175\,\mu\mathrm m$}；高阶反射引起的一阶等效厚度偏移不超过\textbf{$1.2\times10^{-3}\,\mu\mathrm m$}。SiC第三束反射光的最大相对强度仅约$4.85\times10^{-7}\%$，说明所选弱吸收波段内高阶反射对SiC厚度计算的影响很小。

\textbf{本文特点：}构建“\textbf{双光束正演--约束反演--多光束扩展--厚度误差传播}”的递进式建模体系，使SiC与Si在不同吸收条件下的厚度计算具有统一的物理基础；进一步将多光束造成的光谱差异定量传播至厚度尺度，并结合Jacobian--SVD区分模型近似影响与附属参数弱可辨识性。

\keywords{外延层厚度；红外干涉；Sellmeier--Drude模型；Airy多光束模型；非线性最小二乘}